% ============================================================
% Appendix A — Experimental Setup Details and Full Results
% ============================================================

\section{Experimental Setup Details and Full Results}
\label{sec:appendix_setup}

\subsection{Implementation and Execution Details}
\label{sec:appendix_impl}

\paragraph{CSMAD configuration.}
Table~\ref{tab:csmad_config} lists the complete configuration; all values are shared across the
113 entities, and only the input dimensionality $F$ (Table~\ref{tab:dimensionality}), the
data-derived class-prior weight $w_+$ (Eq.~\ref{eq:lcls_app}), and the train/test split
proportions determined by the protocol of Section~\ref{sec:datasets} vary by entity.

%% PDF QA r1 BLOCKER-B: two-column table* + max-width cap; \allowbreak added in
%% the GRL-head row so the layer chain can wrap inside the p{} column.
\begin{table*}[tp]
\footnotesize
\centering
\caption{Complete CSMAD configuration used in all main experiments (single configuration; no
  per-dataset tuning).}
\label{tab:csmad_config}
\setlength{\tabcolsep}{4pt}
\begin{adjustbox}{max width=\linewidth}
\begin{tabular}{llp{8.5cm}}
\toprule
Group & Parameter & Value \\
\midrule
Architecture & Patch embedding & Linear (flatten + projection) \\
             & Encoder         & 4 layers, pre-LN self-attention, $d_{\mathrm{model}}=512$, 8 heads, ff 2048, GELU, dropout 0.15 \\
             & Teacher/Student decoder & 3/2 layers (self-attention only, MAE-style), separate learnable mask tokens \\
             & GRL head & 2-layer MLP: LayerNorm$\to$\allowbreak Linear(512$\to$256)$\to$\allowbreak GELU$\to$\allowbreak Dropout(0.1)$\to$\allowbreak Linear(256$\to$1) \\
\midrule
Windowing & Window/patch size & $L=500$, $s=10$ ($N=50$ patches) \\
          & Train/test stride & 21\,/\,49 \\
          & Masking ratio & 0.15 (8 masked/42 visible); anomaly-priority during training, leave-one-out at test \\
\midrule
Training & Epochs/warmup & 500/first 250 epochs \\
         & Optimizer & AdamW (fused), lr $10^{-3}$, wd $10^{-3}$, $\beta=(0.9,0.99)$ \\
         & GRL lr & $10^{-4}$ (separate parameter group) \\
         & LR schedule & 10-epoch linear warmup + cosine annealing \\
         & Batch/precision/seed & 1024/bf16 mixed/42 (single run per entity) \\
\midrule
Loss & Base weights & $L_{\mathrm{recon}}$, $L_{\mathrm{OD}}$: 1.0; $\lambda_{\mathrm{FM}}$, $\lambda_{\mathrm{GRL}}$: adaptive (Eq.~\ref{eq:adaptive_weight}), $\beta_{\mathrm{FM}}=1.0$, $\beta_{\mathrm{GRL}}=0.2$ \\
     & GRL focal & $\gamma=2$; pos-weight $w_+$ computed per entity \\
\midrule
Scoring & Combination ratio & $c=4$ (Eq.~\ref{eq:sigma}); GRL head unused at inference \\
\bottomrule
\end{tabular}
\end{adjustbox}
\end{table*}

\paragraph{Training budgets and evaluation cadence.}
Table~\ref{tab:budgets} summarizes the per-group training and evaluation budgets reported in
Section~\ref{sec:impl}.
All groups share the best-epoch selection criterion (PA\%K-AUC F1 on the test split), use no
early stopping, complete their full budget, and are reported at their best evaluated epoch.

\begin{table*}[tp]
\footnotesize
\centering
\caption{Training and evaluation budgets per method group.}
\label{tab:budgets}
\begin{adjustbox}{max width=\linewidth}
\begin{tabular}{lcccc}
\toprule
Method group & Epochs & Eval.\ cadence & Batch size & Best-epoch criterion \\
\midrule
CSMAD                          & 500 & every 5 epochs & 1024 & PA\%K-AUC F1 \\
Unsupervised baselines (22)    & 10  & every epoch & model-specific & PA\%K-AUC F1 \\
Weakly supervised baselines (4)& 50  & every epoch & model-specific & PA\%K-AUC F1 \\
\bottomrule
\end{tabular}
\end{adjustbox}
\end{table*}

\paragraph{Environment and code.}
All experiments run on [GPU model]; % PH:TXT-001 | GPU model used for experiments
code, configurations, and the exact dataset partitions will be released at [URL].
% PH:TXT-002 | code repository URL

\paragraph{Baseline implementations and hyperparameters.}
The 22 unsupervised baselines comprise nine detectors adopted from the protocol study
of~\cite{sarfraz2024quovadis} (five simple detectors: random score, sensor-range deviation, PCA
reconstruction, L2-norm, and nearest-neighbor distance; three lightweight neural detectors: MLP,
MLPMixer, and single-stack Transformer; and a GCN-LSTM detector); six established deep MTSAD
systems (Anomaly Transformer, TranAD, USAD, DAGMM, GDN, OmniAnomaly); and seven recent methods
(TFMAE, NPSR, TimesNet, DCdetector, MEMTO, ModernTCN, CATCH).
For DAGMM we follow the simplified re-implementation of the TranAD
repository~\cite{tuli2022tranad} (github.com/imperial-qore/TranAD), in which the GMM energy
term is omitted; the variant is labeled ``DAGMM (simplified)'' in all tables.
Each baseline retains the hyperparameters of its original implementation or publication preset
(for example, NRdetector keeps its native window size of 100 rather than the 500 used by our
pipeline).
The random-score baseline is averaged over five independent runs (mean $\pm$ std), while all
other methods use a single run.
All baselines receive the identical data partitions through a unified loading layer, and all
metrics, for CSMAD and every baseline, are computed by one shared evaluation routine, eliminating
implementation-dependent metric discrepancies.

% ---- Table A.3 placeholder ----
\begin{table*}[tp]
\footnotesize
\centering
\caption{Hyperparameters of all 26 baselines.
  Each method retains the settings of its original implementation or publication preset;
  deviations from the unified pipeline (window size, epochs, batch size) are listed explicitly.
  DAGMM follows the simplified TranAD-repository re-implementation (GMM energy term omitted).}
\label{tab:baseline_hparams}
\begin{adjustbox}{max width=\linewidth}
\begin{tabular}{lccccl}
\toprule
Method & Window & LR & Batch & Epochs & Key parameters \\
\midrule
\multicolumn{6}{l}{\textit{Simple / lightweight (9 methods)}} \\
Random score   & 500 & --- & --- & --- & 5-run average \\
Sensor range   & 500 & --- & --- & --- & per-feature range \\
PCA recon.     & 500 & --- & --- & --- & 50 components \\
L2-norm        & 500 & --- & --- & --- & window L2 \\
NN-distance    & 500 & --- & --- & --- & 5 neighbors \\
MLP            & 500 & {[X.XX]} & {[X.XX]} & 10 & {[X.XX]} \\
MLPMixer       & 500 & {[X.XX]} & {[X.XX]} & 10 & {[X.XX]} \\
Transformer    & 500 & {[X.XX]} & {[X.XX]} & 10 & {[X.XX]} \\
GCN-LSTM       & 500 & {[X.XX]} & {[X.XX]} & 10 & {[X.XX]} \\
\midrule
\multicolumn{6}{l}{\textit{SOTA legacy (6 methods)}} \\
Anomaly Trans. & 100 & {[X.XX]} & {[X.XX]} & 10 & {[X.XX]} \\
TranAD         & 10  & {[X.XX]} & {[X.XX]} & 10 & {[X.XX]} \\
USAD           & 12  & {[X.XX]} & {[X.XX]} & 10 & {[X.XX]} \\
DAGMM (simpl.) & 500 & {[X.XX]} & {[X.XX]} & 10 & GMM omitted \\
GDN            & 15  & {[X.XX]} & {[X.XX]} & 10 & {[X.XX]} \\
OmniAnomaly    & 100 & {[X.XX]} & {[X.XX]} & 10 & {[X.XX]} \\
\midrule
\multicolumn{6}{l}{\textit{SOTA recent (7 methods)}} \\
TFMAE          & 500 & {[X.XX]} & {[X.XX]} & 10 & {[X.XX]} \\
NPSR           & 500 & {[X.XX]} & {[X.XX]} & 10 & {[X.XX]} \\
TimesNet       & 500 & {[X.XX]} & {[X.XX]} & 10 & {[X.XX]} \\
DCdetector     & 105 & {[X.XX]} & {[X.XX]} & 10 & {[X.XX]} \\
MEMTO          & 500 & {[X.XX]} & {[X.XX]} & 10 & {[X.XX]} \\
ModernTCN      & 500 & {[X.XX]} & {[X.XX]} & 10 & {[X.XX]} \\
CATCH          & 500 & {[X.XX]} & {[X.XX]} & 10 & {[X.XX]} \\
\midrule
\multicolumn{6}{l}{\textit{Weakly supervised (4 methods)}} \\
DeepMIL        & 100 & {[X.XX]} & {[X.XX]} & 50 & {[X.XX]} \\
WETAS          & 100 & {[X.XX]} & {[X.XX]} & 50 & {[X.XX]} \\
TreeMIL        & 100 & {[X.XX]} & {[X.XX]} & 50 & {[X.XX]} \\
NRdetector     & 100 & {[X.XX]} & {[X.XX]} & 50 & {[X.XX]} \\
\bottomrule
\end{tabular}
\end{adjustbox}
\end{table*}

\paragraph{SWaT preprocessing (reproducibility note).}
The SWaT (A1+A2) input dimensionality is 45: 51 raw features minus 6 constant columns
\{P202, P401, P404, P502, P601, P603\} removed during preprocessing.
Loading the raw CSV files without the constant-column filter yields 51 features, so reproductions
should verify this dimensionality explicitly.

\subsection{Formal Definitions of Evaluation Metrics}
\label{sec:appendix_metrics}

\paragraph{Point adjustment and PA\%K.}
Under conventional point adjustment (PA)~\cite{xu2018kpivae,kim2022rigorous}, if any timestep
within a ground-truth anomaly segment is predicted positive, all timesteps of that segment are
counted as detected.
PA\%K~\cite{kim2022rigorous} parameterizes this leniency: a segment qualifies for adjustment
only when strictly more than $K\%$ of its timesteps are predicted positive, so that $K=0$
recovers conventional PA and $K=100$ recovers point-wise scoring.

\paragraph{PA\%K-AUC F1.}
For each $K \in \{0,1,\ldots,100\}$, the PA\%K-adjusted F1 is computed with a per-$K$
re-optimized threshold, following the protocol of~\cite{kim2022rigorous}; the reported value is
the trapezoidal integral of F1 over the 101-point grid, normalized to $[0,1]$.
Integrating over the full tolerance spectrum removes the cherry-picking degree of freedom that a
single $K$ would introduce.

\paragraph{PA\%K-AUC AUC-PR.}
The same $K$-grid and integration, with the integrand replaced by the area under the PA\%K-adjusted
precision--recall curve at each $K$ (obtained by a threshold sweep, hence threshold-free).

\paragraph{VUS-ROC\,/\,VUS-PR~\cite{paparrizos2022vus}.}
The Volume Under the Surface generalizes AUC-ROC\,/\,AUC-PR to a three-dimensional volume by
sweeping both the decision threshold and a temporal tolerance parameter that softens segment
boundaries; we use the authors' official implementation with a tolerance window of 100 timesteps,
after min--max normalization of scores.
Both measures are threshold-free and robust to label-boundary uncertainty.

\paragraph{Affiliation F1~\cite{huet2022affiliation}.}
Affiliation precision and recall convert temporal proximity between predicted and ground-truth
events into per-event affinity scores within each event's affiliation zone, with formal
robustness guarantees against adversarial scoring; Affiliation F1 is the harmonic mean of
these two measures.
Binarization uses the anomaly-ratio threshold defined below.

\paragraph{Anomaly-ratio threshold (formal).}
Let $\alpha$ be the fraction of anomalous timesteps in the evaluation span; the threshold is the
$(1-\alpha)$ quantile of the point-level score distribution, and predictions are scores strictly
above it.
The threshold is computed post hoc for threshold-dependent metrics only; it never enters
training or model selection (Section~\ref{sec:impl}).

\paragraph{PA F1 (oracle).}
Conventional PA ($K{=}0$) F1, computed by selecting the F1-optimal threshold on the unadjusted
(pre-PA) predictions and then applying the PA adjustment; reported in \ref{sec:appendix_full_results}
solely for comparability with prior work, marked (oracle), and excluded from all rankings
(Section~\ref{sec:metrics}).

\paragraph{Score aggregation.}
All metrics are computed on point-level scores obtained by mean-aggregation over all covering
(window, patch) pairs (Eq.~\ref{eq:agg}); the identical evaluation routine serves CSMAD and
every baseline.

\subsection{Dataset Details}
\label{sec:appendix_dataset}

\paragraph{Per-entity statistics.}
Table~\ref{tab:per_entity} expands the family-level summary of Table~\ref{tab:datasets}.

\begin{table*}[tp]
\footnotesize
\centering
\caption{Dataset statistics under the contaminated benchmark protocol, per entity.
  Train/test sizes reflect the re-split of Section~\ref{sec:datasets}; Train AR\,/\,Test AR
  denote the anomaly ratio of the training\,/\,evaluation portion.
  SMAP and MSL sizes are concatenated per-channel totals.
  \textit{SMD per-machine rows pending.}}
\label{tab:per_entity}
\begin{adjustbox}{max width=\linewidth}
\begin{tabular}{lrrrccl}
\toprule
Entity & \#Train pts & \#Test pts & \#Dim. & Train AR (\%) & Test AR (\%) & Source \\
\midrule
SWaT (A1+A2) & 719{,}959 & 224{,}960 & 45 & 1.63 & 19.05 (full)\,/\,3.68 & \cite{goh2016swat} \\
WaDi A1      & 1{,}296{,}001 & 86{,}401 & 123 & 0.52 & 3.82 & \cite{ahmed2017wadi} \\
WaDi A2      & 870{,}972 & 86{,}402 & 123 & 0.76 & 3.87 & \cite{ahmed2017wadi} \\
PSM          & 176{,}401 & 43{,}921 & 25 & 6.20 & 30.63 & \cite{abdulaal2021psm} \\
SMD ($\times$28) & [per-machine] & [per-machine] & 29--36 & [per-machine] & 4.16 (avg) & \cite{su2019omnianomaly} \\
SMAP ($\times$54) & 355{,}905 & 217{,}925 & 25 & 0.70 & 24.54 & \cite{hundman2018telemanom} \\
MSL ($\times$27)  & 95{,}271 & 36{,}775 & 55 & 1.70 & 16.72 & \cite{hundman2018telemanom} \\
\bottomrule
\end{tabular}
\end{adjustbox}
\end{table*}

\paragraph{Training-label semantics.}
SWaT and WaDi training files record normal operation only; PSM and SMD distribute no
training-label files, and their training portions are treated as normal, consistent with how
these benchmarks are used in prior work~\cite{su2019omnianomaly,abdulaal2021psm}; SMAP and MSL
training labels are set to zero by the loading pipeline (treated as normal).
Labeled anomalies in our training splits therefore originate exclusively from the incorporated
test-stream prefixes.

\paragraph{Split-boundary adjustment (SMAP/MSL).}
A cut position is accepted only if it lies at least ten timesteps from every annotated anomaly
region; when the 50\% midpoint violates this clearance, the cut moves outward to the nearest
admissible position, without a distance bound.
Table~\ref{tab:split_shifts} reports the measured shifts: 4 of 81 channels (all MSL) moved, and
no boundary-straddling anomaly region remains in any channel.
At the concatenated scale, the aggregate absolute shift is 252 timesteps against an MSL
evaluation total of 36,775.

\begin{table*}[tp]
\footnotesize
\centering
\caption{Measured split-point adjustments for the SMAP/MSL channels
  (all other channels: zero shift).}
\label{tab:split_shifts}
\begin{adjustbox}{max width=\linewidth}
\begin{tabular}{lrrrrl}
\toprule
Channel & Test length & Target (50\%) & Actual cut & Shift & Share of ch.\ test length \\
\midrule
SMAP (all 54) & --- & --- & --- & 0 & 0\% \\
MSL D-16 & 2{,}191 & 1{,}095 & 1{,}261 & $+$166 & 7.58\% \\
MSL M-1  & 2{,}277 & 1{,}138 & 1{,}099 & $-$39  & 1.71\% \\
MSL M-2  & 2{,}277 & 1{,}138 & 1{,}099 & $-$39  & 1.71\% \\
MSL S-2  & 1{,}827 & 913 & 921 & $+$8 & 0.44\% \\
MSL (rem.\ 23) & --- & --- & --- & 0 & 0\% \\
\bottomrule
\end{tabular}
\end{adjustbox}
\end{table*}

\paragraph{Boundary-aware windowing.}
The original training portion and the incorporated test prefix are not temporally adjacent;
segment boundaries are registered so that no sliding window crosses non-contiguous data.
The same mechanism guards the excision boundaries of the anomaly-excised condition
(Section~\ref{sec:baselines}).

\subsection{SWaT Evaluation with Region~22 Excluded (excl22)}
\label{sec:appendix_swat}

\paragraph{Region definition and statistics.}
Attack region~22 is the chronologically first anomaly region within the held-out SWaT evaluation
half, spanning positions $[2{,}869, 38{,}769)$ (indexed within the evaluation span) --- 35,900
contiguous timesteps that constitute 83.75\% of all anomalous timesteps and 15.96\% of the
entire evaluation span.
Its identification is deterministic (first region, with a sanity check on its length; no other
supported dataset contains a single region of comparable extent).
Masking it reduces the evaluation anomaly ratio from 19.05\% to 3.68\%, leaving the 13 smaller
and more diverse attack events.

\paragraph{Evaluation-mask implementation.}
The excl22 condition modifies neither training nor scores: an evaluation mask removes the
region's timesteps from the label vector and the event list before metric computation, so
threshold-free surfaces (VUS) and event-based measures (Affiliation) are likewise computed on
the masked span.
The excl22 condition selects its own best epoch under the shared criterion, independently of
the full condition; the identical mask is applied to every baseline.

% ---- Table A.6 placeholder ----
%% PDF QA r1 BLOCKER-C/BLOCKER-3: 12-column dual-condition table exceeded the
%% text/column width in both builds; table* + \footnotesize + tighter \tabcolsep
%% + max-width cap.
\begin{table*}[tp]
\footnotesize
\centering
\caption{SWaT dual-condition results: all five metrics for CSMAD and all baselines under the
  full condition and the excl22 condition (Section~\ref{sec:datasets}).
  Same trained models and identical scores in both conditions; only the evaluation mask differs.
  The excl22 best epoch is selected independently under the shared criterion.}
\label{tab:swat_dual}
\setlength{\tabcolsep}{3pt}
\begin{adjustbox}{max width=\linewidth}
\begin{tabular}{lccccccccccc}
\toprule
& \multicolumn{5}{c}{Full condition} & & \multicolumn{5}{c}{excl22 condition} \\
\cmidrule(lr){2-6}\cmidrule(lr){8-12}
Method & F1 & PR & VUS-PR & VUS-ROC & Aff. & & F1 & PR & VUS-PR & VUS-ROC & Aff. \\
\midrule
CSMAD      & {[X.XX]} & {[X.XX]} & {[X.XX]} & {[X.XX]} & {[X.XX]} & & {[X.XX]} & {[X.XX]} & {[X.XX]} & {[X.XX]} & {[X.XX]} \\
\multicolumn{12}{c}{\textit{(26 baseline rows; cells {[X.XX]} pending experimental queue)}} \\
\bottomrule
\end{tabular}
\end{adjustbox}
\end{table*}

\subsection{Full Multi-Metric Results}
\label{sec:appendix_full_results}

% ---- Table A.7 placeholder ----
\begin{table*}[tp]
\footnotesize
\centering
\caption{Complete multi-metric results for all methods and dataset families: PA\%K-AUC AUC-PR,
  VUS-ROC, Affiliation F1, and PA F1 (oracle threshold; reported for comparability only, never
  used for ranking --- Section~\ref{sec:metrics}).
  PA\%K-AUC F1 and VUS-PR appear in Table~\ref{tab:main_results}.}
\label{tab:full_metrics}
\begin{adjustbox}{max width=\linewidth}
\begin{tabular}{llccccccc}
\toprule
& & SWaT & WaDi A1 & WaDi A2 & PSM & SMD & SMAP & MSL \\
Method & Metric & excl22 & & & & avg & avg & avg \\
\midrule
CSMAD & AUC-PR   & {[X.XX]} & {[X.XX]} & {[X.XX]} & {[X.XX]} & {[X.XX]} & {[X.XX]} & {[X.XX]} \\
      & VUS-ROC  & {[X.XX]} & {[X.XX]} & {[X.XX]} & {[X.XX]} & {[X.XX]} & {[X.XX]} & {[X.XX]} \\
      & Aff.\ F1 & {[X.XX]} & {[X.XX]} & {[X.XX]} & {[X.XX]} & {[X.XX]} & {[X.XX]} & {[X.XX]} \\
      & PA F1 (oracle) & {[X.XX]} & {[X.XX]} & {[X.XX]} & {[X.XX]} & {[X.XX]} & {[X.XX]} & {[X.XX]} \\
\multicolumn{9}{c}{\textit{(26 baseline rows $\times$ 4 metrics; cells {[X.XX]} pending queue)}} \\
\bottomrule
\end{tabular}
\end{adjustbox}
\end{table*}

\subsection{Per-Entity Results}
\label{sec:appendix_entity_results}

% ---- Table A.8 placeholder ----
\begin{table}[h]
\small
\centering
\caption{Per-entity results (PA\%K-AUC F1\,/\,VUS-PR) for SMD (28 machines), SMAP
  (54 channels), and MSL (27 channels).
  Macro-averages over entities equal the corresponding family columns of
  Table~\ref{tab:main_results}.}
\label{tab:per_entity_results}
\begin{tabular}{lcc}
\toprule
Entity & PA\%K-AUC F1 & VUS-PR \\
\midrule
\multicolumn{3}{l}{\textit{SMD (28 machines; cells {[X.XX]} pending)}} \\
SMD-1-1 & {[X.XX]} & {[X.XX]} \\
\multicolumn{3}{c}{$\vdots$ (28 rows)} \\
\midrule
\multicolumn{3}{l}{\textit{SMAP (54 channels; cells {[X.XX]} pending)}} \\
SMAP-A-1 & {[X.XX]} & {[X.XX]} \\
\multicolumn{3}{c}{$\vdots$ (54 rows)} \\
\midrule
\multicolumn{3}{l}{\textit{MSL (27 channels; cells {[X.XX]} pending)}} \\
MSL-C-1 & {[X.XX]} & {[X.XX]} \\
\multicolumn{3}{c}{$\vdots$ (27 rows)} \\
\bottomrule
\end{tabular}
\end{table}
